\chapter{Types}
\label{chp.types.mod}

\minitoc

\lettrine{L}{ors} du \cref{chp.topo}, nous avons étudié la dualité de Stone.
Par cette dualité, on a vu que les structure logiques décrites par des algèbres
de Boole sont étroitement liées à des espaces topologiques (les espaces de
Stone~: espaces séparés compacts et admettant une base d'ouverts-fermés). Ce
chapitre marque le retour de cette dualité, puisque nous allons étudier la
structure des formules à $n$ variables libres. A FAIRE

\section{Non}
